\begin{center}
\setlength{\parskip}{4pt}

{\LARGE \textbf{\color{tycheblue}Agentic Co-Emergence\\[4pt]
\Large Foundational Artefact 001}}

\vspace{.25cm}

{\large Julian Macnamara and Ariadne\\[4pt]}

{\normalsize June 2026}

\vspace{0.5cm}

\noindent
\makebox[\textwidth][c]{%
\begin{minipage}{0.85\textwidth}

\setlength{\parskip}{0.75em}
\setlength{\parindent}{0pt}

\normalsize
\justifying

This document is the living research artefact for Agentic Co-Emergence. It develops the theoretical foundation of a research programme investigating how understanding emerges and evolves through governed interaction between humans and intelligent agents.

\end{minipage}
}

\end{center}

\section*{1. Motivation}

Agentic Co-Emergence begins from the proposition that the central objective of human-agent systems is not merely reasoning, communication, or decision support.

The deeper objective is the cultivation of understanding.

This requires a shift from treating dialogue as an exchange of messages to treating dialogue as a sequence of governed transformations applied to an evolving understanding.

\section*{2. Research Questions}

The initial research questions are:

\begin{enumerate}
\item What is understanding?
\item How does understanding differ from knowledge, information, and belief?
\item What properties must an understanding possess if it is to be treated as a computational object?
\item What transformations can be applied to an understanding?
\item How can changes in understanding be observed or measured?
\item How do multiple understandings co-emerge?
\item How should understandings be governed?
\item How should understanding be represented computationally?
\end{enumerate}

\section*{3. Axioms}

The following axioms are provisional foundations for the theory.

\begin{enumerate}
\item Understanding is capable of change.
\item Understanding exists only with respect to a domain.
\item Understanding is shaped by purpose.
\item Understanding has structure.
\item Transformations of understanding have provenance.
\item Dialogue can transform understanding.
\item Governance constrains legitimate transformation.
\end{enumerate}

\section*{4. Definitions}

\subsection*{Understanding}

Understanding is a structured, evolving, purpose-relative state that organises concepts, relationships, assumptions, evidence, uncertainties, and tensions within a domain, thereby constraining interpretation, inference, and action.

Understanding is not merely the possession of facts. It is the organised capacity to interpret a domain, relate ideas, recognise uncertainty, revise assumptions, and guide possible action.

\subsection*{Dialogue}

Dialogue is a governed process through which participants exchange, challenge, refine, and integrate perspectives.

Within Agentic Co-Emergence, dialogue is treated as a mechanism of transformation rather than merely communication.

\subsection*{Co-Emergence}

Co-Emergence is the process through which multiple participants collectively produce an understanding that cannot be attributed solely to any individual participant.

\section*{5. Computational Model}

The central computational hypothesis is that understanding can be treated as a first-class computational object.

Such an object must possess:

\begin{itemize}
\item identity;
\item structure;
\item state;
\item domain context;
\item history;
\item transformability;
\item perspective;
\item uncertainty;
\item tension;
\item purpose.
\end{itemize}

The runtime should therefore maintain an evolving understanding object rather than merely storing a transcript of dialogue.

\section*{6. Runtime Architecture}

The initial runtime should provide a minimal representation of an understanding object and the transformations that may be applied to it.

At this stage, the runtime exists to test the theory rather than to provide a complete product.

\section*{7. Experimental Methodology}

Experiments should test whether representing understanding explicitly produces observable benefits.

Initial experiments should investigate whether the system can:

\begin{itemize}
\item represent an understanding state;
\item record changes to that state;
\item distinguish between information, knowledge, and understanding;
\item preserve uncertainty and tension;
\item provide a traceable history of transformation.
\end{itemize}

\section*{8. Discussion}

The early theory suggests that Agentic Co-Emergence is not primarily a theory of artificial intelligence.

It is a theory of how understanding evolves through governed interaction.

Human understanding, agentic understanding, organisational understanding, and shared understanding may therefore be treated as different substrates of a more general phenomenon.

\section*{9. Future Work}

The next stage is to implement the minimal understanding object in the runtime and design tests that challenge whether the representation is meaningful.

The ontology, runtime, experiments, and Living Paper should evolve together.